% !TEX program = xelatex
% ============================================================================
% Diccionario de datos
% Proyecto Final de Bases de Datos - Inventario, Compras y Ventas (PostgreSQL)
% ============================================================================
\documentclass[11pt,oneside]{article}

\usepackage{estilo-proyecto}
\docnombre{Diccionario de datos}

\addto\captionsspanish{\renewcommand{\contentsname}{Contenido}}
\setcounter{tocdepth}{1}

\begin{document}

\portadilla{Proyecto final \textperiodcentered{} Bases de datos}
  {Diccionario de datos}
  {Sistema de inventario, compras y ventas \textperiodcentered{} PostgreSQL}

\tableofcontents
\vspace{0.5em}

\section{Introducci\'on}

Este diccionario documenta las veintid\'os tablas del sistema de gesti\'on de
\textbf{inventario, compras y ventas} implementado en PostgreSQL. La base de
datos adopta una arquitectura de \textit{microservicios por esquema}: cada
servicio es propietario de un esquema independiente y expone sus datos al
resto mediante referencias l\'ogicas.

\begin{center}
\renewcommand{\arraystretch}{1.3}
\begin{tabular}{@{} >{\ttfamily\small}l
                    >{\small}l
                    >{\centering\small\arraybackslash}p{1.5cm} @{}}
\toprule
\rowcolor{fondo}
{\sffamily\bfseries\footnotesize Esquema} &
{\sffamily\bfseries\footnotesize Responsabilidad} &
{\sffamily\bfseries\footnotesize Tablas} \\
\midrule
catalogo    & Maestros: sedes, almacenes, productos y terceros & 7 \\
compras     & \'Ordenes y recepciones de abastecimiento        & 4 \\
inventarios & Existencias y transferencias entre almacenes     & 3 \\
ventas      & Pedidos, pagos y devoluciones                    & 5 \\
identidad   & Empleados, usuarios y auditor\'ia de accesos     & 3 \\
\midrule
            & \sffamily\footnotesize Total                     & 22 \\
\bottomrule
\end{tabular}
\end{center}

\vspace{0.3em}
En las tablas se emplean las siguientes convenciones:

\begin{itemize}[leftmargin=1.5em, itemsep=2pt]
  \item \pk{atributo}: clave primaria, simple o parte de una clave compuesta.
  \item \fk{atributo}: clave for\'anea hacia otra tabla.
  \item \uq{atributo}: restricci\'on de unicidad adicional.
  \item Cuando la clave primaria es compuesta se indica en una nota bajo la tabla.
\end{itemize}

% ============================================================================
\esquema{catalogo}{Cat\'alogo y terceros. Datos maestros del sistema: sedes,
almacenes, productos y actores externos.}

\begin{tabladic}{sedes}
\begin{tabularx}{\linewidth}{A T R}
\toprule
\cabecera
\pk{id\_sede} & INT & Identificador \'unico autoincremental. \\
nombre & VARCHAR(50) & NOT NULL \textperiodcentered{} UNIQUE. Nombre visible de la sede. \\
ciudad & VARCHAR(50) & NOT NULL. Ciudad donde opera la sede. \\
direccion & VARCHAR(80) & NOT NULL. Direcci\'on f\'isica completa. \\
tipo\_sede & VARCHAR(20) & NOT NULL \textperiodcentered{} CHECK IN
  (\texttt{'Ventas'}, \texttt{'Distribuci\'on'}, \texttt{'Regional'}). \\
estado & VARCHAR(10) & NOT NULL \textperiodcentered{} CHECK IN
  (\texttt{'Activo'}, \texttt{'Inactivo'}) \textperiodcentered{} DEFAULT \texttt{'Activo'}. \\
\bottomrule
\end{tabularx}
\end{tabladic}

\begin{tabladic}{almacenes}
\begin{tabularx}{\linewidth}{A T R}
\toprule
\cabecera
\pk{id\_almacen} & INT & Identificador \'unico autoincremental. \\
\fk{id\_sede} & INT & NOT NULL \textperiodcentered{} FK $\to$ \texttt{catalogo.sedes}.
  Sede propietaria del almac\'en. \\
nombre & VARCHAR(50) & NOT NULL. Nombre del almac\'en o bodega. \\
tipo\_almacen & VARCHAR(20) & NOT NULL \textperiodcentered{} CHECK IN
  (\texttt{'Principal'}, \texttt{'Tienda'}, \texttt{'Bodega'}). \\
estado & VARCHAR(10) & NOT NULL \textperiodcentered{} CHECK IN
  (\texttt{'Activo'}, \texttt{'Inactivo'}) \textperiodcentered{} DEFAULT \texttt{'Activo'}. \\
\bottomrule
\end{tabularx}
\end{tabladic}

\begin{tabladic}{categorias}
\begin{tabularx}{\linewidth}{A T R}
\toprule
\cabecera
\pk{id\_categoria} & INT & Identificador \'unico autoincremental. \\
\uq{nombre} & VARCHAR(50) & NOT NULL \textperiodcentered{} UNIQUE. Nombre de la categor\'ia. \\
descripcion & VARCHAR & Descripci\'on libre. Admite nulos. \\
estado & VARCHAR(10) & NOT NULL \textperiodcentered{} CHECK IN
  (\texttt{'Activo'}, \texttt{'Inactivo'}) \textperiodcentered{} DEFAULT \texttt{'Activo'}. \\
\bottomrule
\end{tabularx}
\end{tabladic}

\begin{tabladic}{productos}
\begin{tabularx}{\linewidth}{A T R}
\toprule
\cabecera
\pk{id\_producto} & INT & Identificador \'unico autoincremental. \\
\uq{sku} & VARCHAR(50) & NOT NULL \textperiodcentered{} UNIQUE. C\'odigo de referencia de bodega. \\
nombre & VARCHAR(60) & NOT NULL. Nombre comercial del producto. \\
\fk{id\_categoria} & INT & NOT NULL \textperiodcentered{} FK $\to$ \texttt{catalogo.categorias}. \\
unidad\_medida & VARCHAR(20) & NOT NULL. Por ejemplo: \texttt{'Unidad'}, \texttt{'Kg'}, \texttt{'Litro'}. \\
precio\_lista & NUMERIC(12,2) & NOT NULL \textperiodcentered{} CHECK $\geq 0$.
  Precio de venta base. \\
estado & VARCHAR(10) & NOT NULL \textperiodcentered{} CHECK IN
  (\texttt{'Activo'}, \texttt{'Inactivo'}) \textperiodcentered{} DEFAULT \texttt{'Activo'}. \\
\bottomrule
\end{tabularx}
\end{tabladic}

\begin{tabladic}{proveedores}
\begin{tabularx}{\linewidth}{A T R}
\toprule
\cabecera
\pk{id\_proveedor} & INT & Identificador \'unico autoincremental. \\
\uq{identificacion} & INT & NOT NULL \textperiodcentered{} UNIQUE. NIT o c\'edula del proveedor. \\
razon\_social & VARCHAR(50) & NOT NULL. Nombre legal o raz\'on social. \\
ciudad & VARCHAR(50) & NOT NULL. Ciudad de operaci\'on. \\
correo & VARCHAR(50) & NOT NULL. Correo electr\'onico de contacto. \\
telefono & INT & NOT NULL. N\'umero de contacto. \\
estado & VARCHAR(10) & NOT NULL \textperiodcentered{} CHECK IN
  (\texttt{'Activo'}, \texttt{'Inactivo'}) \textperiodcentered{} DEFAULT \texttt{'Activo'}. \\
\bottomrule
\end{tabularx}
\end{tabladic}

\begin{tabladic}{clientes}
\begin{tabularx}{\linewidth}{A T R}
\toprule
\cabecera
\pk{id\_cliente} & INT & Identificador \'unico autoincremental. \\
tipo\_documento & VARCHAR(15) & NOT NULL \textperiodcentered{} CHECK IN
  (\texttt{'CC'}, \texttt{'NIT'}, \texttt{'CE'}, \texttt{'Pasaporte'}). \\
\uq{num\_documento} & INT & NOT NULL \textperiodcentered{} UNIQUE. N\'umero de documento. \\
nombre\_completo & VARCHAR(60) & NOT NULL. \\
correo & VARCHAR(50) & NOT NULL. \\
telefono & INT & NOT NULL. \\
estado & VARCHAR(10) & NOT NULL \textperiodcentered{} CHECK IN
  (\texttt{'Activo'}, \texttt{'Inactivo'}) \textperiodcentered{} DEFAULT \texttt{'Activo'}. \\
\bottomrule
\end{tabularx}
\end{tabladic}

\begin{tabladic}{producto\_proveedor}
\begin{tabularx}{\linewidth}{A T R}
\toprule
\cabecera
\fk{id\_producto} & INT & NOT NULL \textperiodcentered{} FK $\to$ \texttt{catalogo.productos}.
  \textcolor{pk}{\textbf{Parte de la clave primaria.}} \\
\fk{id\_proveedor} & INT & NOT NULL \textperiodcentered{} FK $\to$ \texttt{catalogo.proveedores}.
  \textcolor{pk}{\textbf{Parte de la clave primaria.}} \\
costo\_referencia & NUMERIC(12,2) & NOT NULL \textperiodcentered{} CHECK $\geq 0$.
  Costo de compra acordado con el proveedor. \\
plazo\_entrega & INT & NOT NULL \textperiodcentered{} CHECK $\geq 0$.
  D\'ias de entrega pactados. \\
\bottomrule
\end{tabularx}
\nota{Clave primaria compuesta: (\texttt{id\_producto}, \texttt{id\_proveedor}).
Tabla asociativa de la relaci\'on M:N \textit{suministra}.}
\end{tabladic}

% ============================================================================
\esquema{compras}{Compras y abastecimiento. Ciclo de aprovisionamiento desde
la orden hasta la recepci\'on f\'isica.}

\begin{tabladic}{ordenes\_compra}
\begin{tabularx}{\linewidth}{A T R}
\toprule
\cabecera
\pk{id\_orden\_compra} & INT & Identificador \'unico autoincremental. \\
\fk{id\_proveedor} & INT & NOT NULL \textperiodcentered{} FK $\to$ \texttt{catalogo.proveedores}
  (referencia l\'ogica entre esquemas). \\
\fk{id\_almacen\_destino} & INT & NOT NULL \textperiodcentered{} FK $\to$ \texttt{catalogo.almacenes}.
  Bodega receptora. \\
\fk{id\_empleado} & INT & NOT NULL \textperiodcentered{} FK $\to$ \texttt{identidad.empleados}.
  Responsable de la orden. \\
fecha\_orden & DATE & NOT NULL. Fecha de emisi\'on. \\
fecha\_estimada & DATE & NOT NULL \textperiodcentered{} CHECK $\geq$ \texttt{fecha\_orden}.
  Entrega pactada. \\
estado & VARCHAR(20) & NOT NULL \textperiodcentered{} CHECK IN
  (\texttt{'Pendiente'}, \texttt{'Recibida Parcial'}, \texttt{'Recibida Total'},
  \texttt{'Cancelada'}). \\
\bottomrule
\end{tabularx}
\end{tabladic}

\begin{tabladic}{detalle\_orden\_compra}
\begin{tabularx}{\linewidth}{A T R}
\toprule
\cabecera
\fk{id\_orden\_compra} & INT & NOT NULL \textperiodcentered{} FK $\to$ \texttt{compras.ordenes\_compra}.
  \textcolor{pk}{\textbf{Parte de la clave primaria.}} \\
\fk{id\_producto} & INT & NOT NULL \textperiodcentered{} FK $\to$ \texttt{catalogo.productos}.
  \textcolor{pk}{\textbf{Parte de la clave primaria.}} \\
cantidad\_solicitada & INT & NOT NULL \textperiodcentered{} CHECK $> 0$. \\
costo\_unitario & NUMERIC(12,2) & NOT NULL \textperiodcentered{} CHECK $\geq 0$.
  Precio negociado en la orden. \\
\bottomrule
\end{tabularx}
\nota{Clave primaria compuesta: (\texttt{id\_orden\_compra}, \texttt{id\_producto}).
Implementa la agregaci\'on A1 del modelo entidad-relaci\'on.}
\end{tabladic}

\begin{tabladic}{recepciones\_compra}
\begin{tabularx}{\linewidth}{A T R}
\toprule
\cabecera
\pk{id\_recepcion} & INT & Identificador \'unico autoincremental. \\
\fk{id\_orden\_compra} & INT & NOT NULL \textperiodcentered{} FK $\to$ \texttt{compras.ordenes\_compra}. \\
fecha\_recepcion & DATE & NOT NULL. Fecha de ingreso f\'isico a bodega. \\
estado & VARCHAR(20) & NOT NULL \textperiodcentered{} CHECK IN
  (\texttt{'Procesado'}, \texttt{'Parcial'}, \texttt{'Anulado'}). \\
\bottomrule
\end{tabularx}
\end{tabladic}

\begin{tabladic}{detalle\_recepcion}
\begin{tabularx}{\linewidth}{A T R}
\toprule
\cabecera
\fk{id\_recepcion} & INT & NOT NULL \textperiodcentered{} FK $\to$ \texttt{compras.recepciones\_compra}.
  \textcolor{pk}{\textbf{Parte de la clave primaria.}} \\
\fk{id\_producto} & INT & NOT NULL \textperiodcentered{} FK $\to$ \texttt{catalogo.productos}.
  \textcolor{pk}{\textbf{Parte de la clave primaria.}} \\
cantidad\_recibida & INT & NOT NULL \textperiodcentered{} CHECK $\geq 0$.
  Unidades f\'isicamente ingresadas. \\
\bottomrule
\end{tabularx}
\nota{Clave primaria compuesta: (\texttt{id\_recepcion}, \texttt{id\_producto}).}
\end{tabladic}

% ============================================================================
\esquema{inventarios}{Inventarios y transferencias. Existencias por almac\'en
y movimientos de mercanc\'ia entre sedes.}

\begin{tabladic}{inventarios}
\begin{tabularx}{\linewidth}{A T R}
\toprule
\cabecera
\fk{id\_sede} & INT & NOT NULL \textperiodcentered{} FK $\to$ \texttt{catalogo.sedes}.
  \textit{Clave de partici\'on LIST.} \\
\fk{id\_almacen} & INT & NOT NULL \textperiodcentered{} FK $\to$ \texttt{catalogo.almacenes}.
  \textcolor{pk}{\textbf{Parte de la clave primaria.}} \\
\fk{id\_producto} & INT & NOT NULL \textperiodcentered{} FK $\to$ \texttt{catalogo.productos}.
  \textcolor{pk}{\textbf{Parte de la clave primaria.}} \\
existencia\_actual & INT & NOT NULL \textperiodcentered{} CHECK $\geq 0$.
  Unidades disponibles en el almac\'en. \\
stock\_minimo & INT & NOT NULL \textperiodcentered{} CHECK $\geq 0$.
  Umbral de alerta de reabastecimiento. \\
fecha\_actualizacion & DATE & NOT NULL. \'Ultima fecha de movimiento registrado. \\
\bottomrule
\end{tabularx}
\nota{Clave primaria compuesta: (\texttt{id\_almacen}, \texttt{id\_producto}).
Tabla particionada por \texttt{id\_sede} (LIST). Implementa la agregaci\'on A2.}
\end{tabladic}

\begin{tabladic}{transferencias}
\begin{tabularx}{\linewidth}{A T R}
\toprule
\cabecera
\pk{id\_transferencia} & INT & Identificador \'unico autoincremental. \\
\fk{id\_almacen\_origen} & INT & NOT NULL \textperiodcentered{} FK $\to$ \texttt{catalogo.almacenes}. \\
\fk{id\_almacen\_destino} & INT & NOT NULL \textperiodcentered{} FK $\to$ \texttt{catalogo.almacenes}. \\
fecha\_solicitud & DATE & NOT NULL. Fecha de solicitud del traslado. \\
fecha\_despacho & DATE & Admite nulos. Fecha de salida del almac\'en origen. \\
fecha\_recepcion & DATE & Admite nulos. Fecha de llegada al almac\'en destino. \\
estado & VARCHAR(20) & NOT NULL \textperiodcentered{} CHECK IN
  (\texttt{'Solicitada'}, \texttt{'Despachada'}, \texttt{'Recibida'}, \texttt{'Cancelada'}). \\
\bottomrule
\end{tabularx}
\nota{Restricciones adicionales: \texttt{id\_almacen\_origen <> id\_almacen\_destino};
\texttt{fecha\_despacho >= fecha\_solicitud}; \texttt{fecha\_recepcion >= fecha\_despacho}.}
\end{tabladic}

\begin{tabladic}{detalle\_transferencia}
\begin{tabularx}{\linewidth}{A T R}
\toprule
\cabecera
\fk{id\_transferencia} & INT & NOT NULL \textperiodcentered{} FK $\to$ \texttt{inventarios.transferencias}.
  \textcolor{pk}{\textbf{Parte de la clave primaria.}} \\
\fk{id\_producto} & INT & NOT NULL \textperiodcentered{} FK $\to$ \texttt{catalogo.productos}.
  \textcolor{pk}{\textbf{Parte de la clave primaria.}} \\
cantidad\_enviada & INT & NOT NULL \textperiodcentered{} CHECK $> 0$. Unidades despachadas. \\
cantidad\_recibida & INT & CHECK $\geq 0$ o nulo. Unidades confirmadas en destino. \\
\bottomrule
\end{tabularx}
\nota{Clave primaria compuesta: (\texttt{id\_transferencia}, \texttt{id\_producto}).}
\end{tabladic}

% ============================================================================
\esquema{ventas}{Ventas y pagos. Pedidos, medios de pago y devoluciones de
mercanc\'ia.}

\begin{tabladic}{pedidos\_venta}
\begin{tabularx}{\linewidth}{A T R}
\toprule
\cabecera
\pk{id\_pedido} & INT & Identificador \'unico autoincremental. \\
\fk{id\_sede} & INT & NOT NULL \textperiodcentered{} FK $\to$ \texttt{catalogo.sedes}.
  \textit{Clave de partici\'on LIST.} \\
\fk{id\_almacen} & INT & NOT NULL \textperiodcentered{} FK $\to$ \texttt{catalogo.almacenes}.
  Almac\'en de despacho. \\
\fk{id\_cliente} & INT & NOT NULL \textperiodcentered{} FK $\to$ \texttt{catalogo.clientes}. \\
\fk{id\_empleado} & INT & NOT NULL \textperiodcentered{} FK $\to$ \texttt{identidad.empleados}.
  Vendedor responsable. \\
fecha\_pedido & DATE & NOT NULL. \\
estado & VARCHAR(20) & NOT NULL \textperiodcentered{} CHECK IN
  (\texttt{'Pendiente'}, \texttt{'Completado'}, \texttt{'Cancelado'}). \\
total\_registrado & NUMERIC(12,2) & NOT NULL \textperiodcentered{} CHECK $\geq 0$.
  Total calculado del pedido. \\
\bottomrule
\end{tabularx}
\nota{Tabla particionada por \texttt{id\_sede} (LIST): un fragmento por cada sede de ventas.}
\end{tabladic}

\begin{tabladic}{detalle\_pedido}
\begin{tabularx}{\linewidth}{A T R}
\toprule
\cabecera
\fk{id\_pedido} & INT & NOT NULL \textperiodcentered{} FK $\to$ \texttt{ventas.pedidos\_venta}.
  \textcolor{pk}{\textbf{Parte de la clave primaria.}} \\
\fk{id\_producto} & INT & NOT NULL \textperiodcentered{} FK $\to$ \texttt{catalogo.productos}.
  \textcolor{pk}{\textbf{Parte de la clave primaria.}} \\
cantidad & INT & NOT NULL \textperiodcentered{} CHECK $> 0$. \\
precio\_unitario & NUMERIC(12,2) & NOT NULL \textperiodcentered{} CHECK $\geq 0$.
  Precio real de venta en el momento del pedido. \\
\bottomrule
\end{tabularx}
\nota{Clave primaria compuesta: (\texttt{id\_pedido}, \texttt{id\_producto}).
Implementa la agregaci\'on A3.}
\end{tabladic}

\begin{tabladic}{pagos}
\begin{tabularx}{\linewidth}{A T R}
\toprule
\cabecera
\pk{id\_pago} & INT & Identificador \'unico autoincremental. \\
\fk{id\_pedido} & INT & NOT NULL \textperiodcentered{} FK $\to$ \texttt{ventas.pedidos\_venta}. \\
fecha\_pago & DATE & NOT NULL. \\
valor & NUMERIC(12,2) & NOT NULL \textperiodcentered{} CHECK $> 0$. \\
medio\_pago & VARCHAR(20) & NOT NULL \textperiodcentered{} CHECK IN
  (\texttt{'Efectivo'}, \texttt{'Tarjeta Debito'}, \texttt{'Tarjeta Credito'},
  \texttt{'Transferencia'}). \\
estado & VARCHAR(20) & NOT NULL \textperiodcentered{} CHECK IN
  (\texttt{'Confirmado'}, \texttt{'Rechazado'}, \texttt{'Pendiente'}). \\
\bottomrule
\end{tabularx}
\end{tabladic}

\begin{tabladic}{devoluciones}
\begin{tabularx}{\linewidth}{A T R}
\toprule
\cabecera
\pk{id\_devolucion} & INT & Identificador \'unico autoincremental. \\
\fk{id\_pedido} & INT & NOT NULL \textperiodcentered{} FK $\to$ \texttt{ventas.pedidos\_venta}. \\
fecha\_solicitud & DATE & NOT NULL. Fecha de radicaci\'on de la devoluci\'on. \\
motivo & VARCHAR & NOT NULL. Descripci\'on del motivo. \\
estado & VARCHAR(20) & NOT NULL \textperiodcentered{} CHECK IN
  (\texttt{'Solicitada'}, \texttt{'Aprobada'}, \texttt{'Rechazada'}). \\
\bottomrule
\end{tabularx}
\end{tabladic}

\begin{tabladic}{detalle\_devolucion}
\begin{tabularx}{\linewidth}{A T R}
\toprule
\cabecera
\fk{id\_devolucion} & INT & NOT NULL \textperiodcentered{} FK $\to$ \texttt{ventas.devoluciones}.
  \textcolor{pk}{\textbf{Parte de la clave primaria.}} \\
\fk{id\_producto} & INT & NOT NULL \textperiodcentered{} FK $\to$ \texttt{catalogo.productos}.
  \textcolor{pk}{\textbf{Parte de la clave primaria.}} \\
cantidad & INT & NOT NULL \textperiodcentered{} CHECK $> 0$. Unidades devueltas. \\
condicion\_producto & VARCHAR(20) & NOT NULL \textperiodcentered{} CHECK IN
  (\texttt{'Buen Estado'}, \texttt{'Defectuoso'}, \texttt{'Da\~nado'}). \\
\bottomrule
\end{tabularx}
\nota{Clave primaria compuesta: (\texttt{id\_devolucion}, \texttt{id\_producto}).}
\end{tabladic}

% ============================================================================
\esquema{identidad}{Identidad y accesos. Personal, credenciales del sistema y
auditor\'ia de operaciones.}

\begin{tabladic}{empleados}
\begin{tabularx}{\linewidth}{A T R}
\toprule
\cabecera
\pk{id\_empleado} & INT & Identificador \'unico autoincremental. \\
\fk{id\_sede} & INT & NOT NULL \textperiodcentered{} FK $\to$ \texttt{catalogo.sedes}.
  Sede asignada al empleado. \\
\uq{num\_documento} & INT & NOT NULL \textperiodcentered{} UNIQUE. C\'edula o documento de identidad. \\
nombre\_completo & VARCHAR(60) & NOT NULL. \\
cargo & VARCHAR(50) & NOT NULL. Por ejemplo: \texttt{'Vendedor'},
  \texttt{'Jefe de Bodega'}, \texttt{'Gerente'}. \\
correo & VARCHAR(50) & NOT NULL. Correo institucional. \\
estado & VARCHAR(10) & NOT NULL \textperiodcentered{} CHECK IN
  (\texttt{'Activo'}, \texttt{'Inactivo'}) \textperiodcentered{} DEFAULT \texttt{'Activo'}. \\
\bottomrule
\end{tabularx}
\end{tabladic}

\begin{tabladic}{usuarios}
\begin{tabularx}{\linewidth}{A T R}
\toprule
\cabecera
\pk{id\_usuario} & INT & Identificador \'unico autoincremental. \\
\fk{\uq{id\_empleado}} & INT & NOT NULL \textperiodcentered{} UNIQUE \textperiodcentered{}
  FK $\to$ \texttt{identidad.empleados}. Relaci\'on 1:1 empleado--usuario. \\
\uq{nombre\_usuario} & VARCHAR(50) & NOT NULL \textperiodcentered{} UNIQUE.
  Nombre de inicio de sesi\'on. \\
password\_hash & VARCHAR(255) & NOT NULL. Hash bcrypt de la contrase\~na. \\
rol & VARCHAR(20) & NOT NULL \textperiodcentered{} CHECK IN
  (\texttt{'Administrador'}, \texttt{'Gerencia'}, \texttt{'Inventario'},
  \texttt{'Compras'}, \texttt{'Vendedor'}, \texttt{'Auditor'}). \\
estado & VARCHAR(10) & NOT NULL \textperiodcentered{} CHECK IN
  (\texttt{'Activo'}, \texttt{'Inactivo'}) \textperiodcentered{} DEFAULT \texttt{'Activo'}. \\
fecha\_creacion & DATE & NOT NULL. Fecha de alta del usuario en el sistema. \\
\bottomrule
\end{tabularx}
\end{tabladic}

\begin{tabladic}{auditoria\_accesos}
\begin{tabularx}{\linewidth}{A T R}
\toprule
\cabecera
\pk{id\_auditoria} & INT & Identificador \'unico autoincremental. \\
\fk{id\_usuario} & INT & NOT NULL \textperiodcentered{} FK $\to$ \texttt{identidad.usuarios}. \\
fecha\_hora & TIMESTAMP & NOT NULL. Momento exacto de la acci\'on registrada. \\
accion & VARCHAR(50) & NOT NULL. Tipo de operaci\'on:
  \texttt{'SELECT'}, \texttt{'INSERT'}, \texttt{'UPDATE'}, \texttt{'DELETE'}. \\
objeto\_consultado & VARCHAR(80) & NOT NULL. Tabla o vista sobre la que se
  ejecut\'o la acci\'on. \\
\bottomrule
\end{tabularx}
\nota{Tabla de solo escritura para los roles de auditor\'ia: sin UPDATE ni DELETE
por pol\'itica de seguridad.}
\end{tabladic}

\end{document}
